\documentclass[10pt,a4paper]{amsart}
\usepackage[margin=19mm]{geometry}
\usepackage{amsmath,amssymb,mathtools}
\usepackage[scr=rsfs,cal=euler]{mathalfa}
\usepackage{xfrac}
\usepackage[unicode,hidelinks,bookmarks=false]{hyperref}
\newcommand{\ie}{i.e.\ }
\newcommand{\cf}{cf.\ }
\newcommand{\resp}{respectively\ }
\newcommand{\Subsection}{\subsection}
\providecommand{\nobreakdash}{\nobreak\hbox{-}}
\setlength{\emergencystretch}{2em}
\multlinegap0pt
\usepackage{bm}
\usepackage{bbm}
\usepackage{dsfont}
\usepackage{xspace}
\usepackage{cancel}
\usepackage{mathtools}
\usepackage{amsthm}
\makeatletter
\@ifundefined{theorem}{\newtheorem{theorem}{Theorem}}{}
\makeatother
\newcommand{\RR}{{\boldsymbol{R}}}
\newcommand{\bd}{\partial}
\newcommand{\lra}{\longrightarrow}
\title[Bounded domains in the 3-dimensional space]{Bounded domains in the 3-dimensional space}
\author{Takashi Tsuboi}
\date{Source: arXiv:2511.01278v1 (CC BY 4.0)}
\begin{document}
\maketitle

\noindent\emph{Source and licence.} Bounded domains in the 3-dimensional space. Version arXiv:2511.01278v1 (CC BY 4.0), distributed under the Creative Commons Attribution 4.0 International licence, as recorded in the arXiv licence metadata for this exact version and shown on the abstract page https://arxiv.org/abs/2511.01278v1. Full rights notice: licenses/arxiv-2511.01278-CC-BY-4.0.txt.\par\smallskip

\noindent\textit{Editorial scope.} Counted unit: the Theorem whose source label is \texttt{th:without\_concave\_critial\_points} in the article (source0.tex lines 923--925, source label \texttt{th:without\_concave\_critial\_points}), delivered together with the article's own complete original proof of it (source0.tex lines 926--963, one \texttt{proof} environment of 38 lines). No mathematics is written, repaired or re-derived: statement and proof are the article's own text line for line. Required context retained uncounted: none. Every definition, display and label of those lines is retained unchanged. Omitted from this package and not redistributed: 1--922, 964--1176. Nothing inside a retained excerpt is omitted or altered: every retained body is delivered exactly as the article's lines are written, so no omission receipt of any kind accompanies this package. Citations: the retained excerpts carry no citation command at all, so no bibliography entry of the article is transcribed and no reference is redistributed. Reference adaptation: none was needed and none was made; every reference inside the delivered unit resolves inside it, so no unresolved reference or citation is produced. Printed slips kept literal: no printed word, symbol, hypothesis or number of the counted statement or of its proof was altered. Layout and class: the article's own class is \texttt{amsart} with packages including graphicx, amsmath, amssymb, amsbsy, latexsym, color, amscd; the wrapper is the lane's pinned \texttt{amsart} editorial wrapper at 10pt on A4, which restates the theorem-like environments the retained text needs and adds the article's own macro definitions that the retained text needs (\texttt{\textbackslash RR}, \texttt{\textbackslash bd}, \texttt{\textbackslash lra}), with extra wrapper packages (bm, bbm, dsfont, xspace, cancel, mathtools). The delivered numbering is the unit's own. Assets: no retained span contains a figure environment, a graphics-inclusion command, a PDF-page inclusion, a table or a TikZ picture; no image file is copied into this package, the package declares no asset dependency and no image file is redistributed with it. Licence: version arXiv:2511.01278v1 of this article is distributed under the Creative Commons Attribution 4.0 International licence, as recorded in the arXiv licence metadata for this exact version and shown on the abstract page https://arxiv.org/abs/2511.01278v1. A full rights notice is delivered at licenses/arxiv-2511.01278-CC-BY-4.0.txt. Characters: every retained excerpt is pure ASCII, so the delivered engine, XeLaTeX with its default Latin Modern text font, reports no missing character.
\begin{theorem}\label{th:without_concave_critial_points}
For a connected boundary connected bounded domain $M$ in $\RR^3$, let $F:M\lra \RR$ be a Morse height function. If there are no concave critical points of index 2 of $F|\bd M$, then $M$ is an embedded handlebody. The same is true if there are no concave critical points of index 0 of $F|\bd M$.
\end{theorem}

\begin{proof}
Assume that the Morse height function $F$ has minimum number of critical points among the isotopy class of the bounded domain $M$ with no concave critical points of index 0.

We look at the minus gradient flow of the Morse function restricted to $\bd M$.
For the minus gradient flow, there are finitely many orbits which does not goes to the local minimum points. They are the local maximum points and the orbits which goes to saddle points. We may assume that the critical points of $F|\bd M$ have distinct (critical) values and there are no orbits connecting two saddles. 

We look at the change in shape of $F^{-1}((-\infty,z])$ as $z$ varies from the minimum value to the maximum value.
At the minimum value $z$ of $F|\bd M$, the critical point $(x,y,z)$ is of index 0 and convex. Hence in a neighborhood of $(x,y,z)$, $F^{-1}((-\infty,z+\varepsilon])$ for small positive $\varepsilon$ is a 3-disk, and it is contractible. 

As $z$ increases, we may have critical points of index 0, critical points of index 1 or convex critical points of index 2. We look at what happens as $z$ passes through the critical values in each case.

When $z$ passes through the critical value of a convex critical point of index 0, a 3-disk appears in $F^{-1}((-\infty,z+\varepsilon])$ 
and a contractible connected component is added. 
When $z$ passes through the critical value of a concave critical point $(x,y,z)$ of index 0, there appears a hole in $F^{-1}((-\infty,z+\varepsilon])$ with the deepest point at the critical point $(x,y,z)$, and the homotopy type of $F^{-1}((-\infty,z+\varepsilon])$ is the same as that of $F^{-1}((-\infty,z-\varepsilon])$.  $F^{-1}((-\infty,z-\varepsilon])$ is a deformation retract of $F^{-1}((-\infty,z+\varepsilon])$.

When $z$ passes through the critical value of a critical point $(x,y,z)$ of index 1, there are two cases; either the outward normal vector at $(x,y,z)$ is upward or it is downward. 

If the outward normal vector at the critical point $(x,y,z)$ is upward, the level set $F^{-1}(z-\varepsilon)$ near $(x,y,z)$ splits into two parts to the level set $F^{-1}(z+\varepsilon)$ near $(x,y,z)$. This does not change the homotopy type, that is, the homotopy type of $F^{-1}((-\infty,z+\varepsilon])$ is the same as that of  $F^{-1}((-\infty,z-\varepsilon])$. $F^{-1}((-\infty,z-\varepsilon])$ is a deformation retract of $F^{-1}((-\infty,z+\varepsilon])$.

If the outward normal vector at the critical point $(x,y,z)$ is downward, the two components of the level set $F^{-1}(z-\varepsilon)$ near $(x,y,z)$ becomes one connected component of the level set $F^{-1}(z+\varepsilon)$. This corresponds to attaching a solid 1-handle to $F^{-1}((-\infty,z-\varepsilon])$. This solid 0ne handle can be seen as a collar neighborhood of the band 1-handle attached to $(F|\bd M)^{-1}(-\infty,z-\varepsilon])$. 
The homotopy type of $F^{-1}((-\infty,z+\varepsilon])$ is that of the union of 
  $F^{-1}((-\infty,z-\varepsilon])$ and the stable manifold (arc) of the critial point $(x,y,z)$ of the gradient flow for $F|\bd M$ on $\bd M$. Thus if  $F^{-1}((-\infty,z-\varepsilon])$ has the homotopy type of 1-dimensional cellar complex then  $F^{-1}((-\infty,z+\varepsilon])$ also has the homotopy type of 1-dimensional cellar complex.  This attaching 1-handle causes either making separated components connected or forming a circle which is non-trivial in the fundamental group. 

When $z$ passes through the critical value of a convex critical point $(x,y,z)$ of index 2, then this is a maximal point and the homotopy type of $F^{-1}((-\infty,z+\varepsilon])$ is the same as that of  $F^{-1}((-\infty,z-\varepsilon])$. $F^{-1}((-\infty,z-\varepsilon])$ is a deformation retract of $F^{-1}((-\infty,z+\varepsilon])$.

These local arguments show that the bounded domain $M$ has the homotopy type of 1 dimensional cellar complex. This implies that $M$ is isotopic to a handlebody or a regular neighborhood of a spacial graph. 

More explicitly, 
the 1-dimensional cellar complex is the union of following parts:
\begin{itemize}
\item 
the stable manifolds (arc) of the critial points  $(x,y,z)$ of index 1 with downward outward normal vector of the gradient flow for $F|\bd M$ on $\bd M$.
\item
when such a stable manifolds is an unstable arc of a concave critical point $p_1$ of index 0, the vertical line segment $p_2p_1$ in $M$ with $p_1$, $p_2\in \bd M$, and the flow line of the gradient flow for $F|\bd M$ from a convex critial point $p_3$ of index 0 to $p_2$. Here we note that generically this flow line is not from a critial point of index 1.
\end{itemize}

This 1-dimensional cellar complex is a deformationretract of $M$ and $M$ is isotopic to a regular neighborhood of this 1-dimensional cellar complex.
\end{proof}

\end{document}
